\chapter{Two-Dimensional Conformal Symmetry and the Virasoro Algebra}
\label{ch:volume-iii-two-dimensional-virasoro}

\section{Complex Coordinates}

Let a two-dimensional Euclidean CFT be placed on the complex plane with
\[
  z=x^1+\ii x^2,
  \qquad
  \bar z=x^1-\ii x^2.
\]
The flat metric is
\[
  \dd s^2=\dd z\,\dd\bar z .
\]
The local conformal transformations are
\[
  z\mapsto f(z),
  \qquad
  \bar z\mapsto \bar f(\bar z),
\]
where \(f\) is holomorphic in the coordinate patch.  Infinitesimally,
\[
  \delta z=\epsilon(z),
  \qquad
  \delta\bar z=\bar\epsilon(\bar z).
\]

The stress tensor decomposes into holomorphic and antiholomorphic parts:
\[
  T(z)=T_{zz}(z),
  \qquad
  \widetilde T(\bar z)=T_{\bar z\bar z}(\bar z).
\]
Conservation and tracelessness imply
\[
  \bar\partial T(z)=0,
  \qquad
  \partial \widetilde T(\bar z)=0
\]
away from insertions.  Thus \(T\) and \(\widetilde T\) are chiral local
fields.

\section{Virasoro Modes}

For a closed contour \(C\) around the origin, define
\begin{equation}
  L_n
  =
  \oint_C\frac{\dd z}{2\pi\ii}\,z^{n+1}T(z),
  \qquad
  \widetilde L_n
  =
  \oint_C\frac{\dd\bar z}{2\pi\ii}\,
  \bar z^{n+1}\widetilde T(\bar z).
  \label{eq:virasoro-mode-definition}
\end{equation}
The contour form makes the generator local: moving the contour through a
correlation function records the residues of the stress-tensor OPEs with the
insertions it crosses.

\begin{figure}[H]
  \centering
  \resizebox{0.86\textwidth}{!}{%
  \begin{tikzpicture}[x=1cm,y=1cm,every node/.style={font=\scriptsize}]
    \tikzset{
      insertion/.style={circle,fill=violet!80!black,inner sep=1.8pt},
      arrow/.style={-{Latex[length=2mm]},line width=0.8pt},
      contour/.style={green!45!black,line width=0.85pt}
    }
    \begin{scope}
      \draw[->] (-2.0,0)--(2.0,0) node[right] {\(\Re z\)};
      \draw[->] (0,-1.65)--(0,1.75) node[above] {\(\Im z\)};
      \draw[contour] (0,0) circle (0.95);
      \node[insertion] at (0,0) {};
      \node[violet!80!black,above right] at (0,0) {\(\mathcal O(0)\)};
      \node[green!45!black] at (-0.92,1.05) {\(C\)};
      \node at (0,-2.05)
      {\(\displaystyle L_n=\oint_C\frac{\dd z}{2\pi\ii}z^{n+1}T(z)\)};
    \end{scope}

    \draw[arrow] (2.65,0)--(3.90,0);
    \node[align=center] at (3.28,0.50) {move\\contour};

    \begin{scope}[shift={(6.15,0)}]
      \draw[contour] (0,0) circle (1.15);
      \draw[contour,orange!85!black] (0,0) circle (0.65);
      \node[insertion] at (0,0) {};
      \node[violet!80!black,above right] at (0,0) {\(\phi(0)\)};
      \node[green!45!black] at (-1.05,1.10) {\(C_1\)};
      \node[orange!85!black] at (0.85,0.72) {\(C_2\)};
      \node[blue!65!black] at (0,-1.55)
      {commutators are residues};
    \end{scope}
  \end{tikzpicture}
  }
  \caption{Virasoro generators are contour integrals of the chiral stress
  tensor.  Deforming contours through insertions computes their commutators.}
  \label{fig:virasoro-contours}
\end{figure}

\section{Primary Fields}

A local field \(\phi(z,\bar z)\) has weights \((h,\bar h)\) if its scaling
dimension and spin are
\[
  \Delta=h+\bar h,
  \qquad
  s=h-\bar h.
\]
It is a Virasoro primary when its stress-tensor OPEs are
\begin{align}
  T(z)\phi(w,\bar w)
  &\sim
  \frac{h\,\phi(w,\bar w)}{(z-w)^2}
  +
  \frac{\partial_w\phi(w,\bar w)}{z-w},
  \label{eq:T-primary-ope}\\
  \widetilde T(\bar z)\phi(w,\bar w)
  &\sim
  \frac{\bar h\,\phi(w,\bar w)}{(\bar z-\bar w)^2}
  +
  \frac{\partial_{\bar w}\phi(w,\bar w)}{\bar z-\bar w}.
  \label{eq:Tbar-primary-ope}
\end{align}
Equivalently, under \(z\mapsto f(z)\),
\[
  \phi'(f(z),\bar f(\bar z))
  =
  \left(\frac{\dd f}{\dd z}\right)^{-h}
  \left(\frac{\dd\bar f}{\dd\bar z}\right)^{-\bar h}
  \phi(z,\bar z).
\]
The mode commutator follows by integrating \eqref{eq:T-primary-ope}:
\begin{equation}
  [L_n,\phi(z,\bar z)]
  =
  \left(z^{n+1}\partial_z+h(n+1)z^n\right)\phi(z,\bar z).
  \label{eq:Ln-primary-commutator}
\end{equation}
At the origin, the state \(\ket{\phi}\) associated with a primary obeys
\[
  L_0\ket{\phi}=h\ket{\phi},
  \qquad
  L_n\ket{\phi}=0\quad(n>0),
\]
and similarly for \(\widetilde L_n\).

\section{The Stress-Tensor OPE}

The chiral stress tensor has the self-OPE
\begin{equation}
  T(z)T(w)
  \sim
  \frac{c/2}{(z-w)^4}
  +
  \frac{2T(w)}{(z-w)^2}
  +
  \frac{\partial T(w)}{z-w}.
  \label{eq:TT-ope}
\end{equation}
The number \(c\) is the central charge.  Inserting
\eqref{eq:TT-ope} into the double contour integral defining
\([L_n,L_m]\) gives
\begin{equation}
  [L_n,L_m]
  =
  (n-m)L_{n+m}
  +
  \frac{c}{12}(n^3-n)\delta_{n,-m}.
  \label{eq:virasoro-algebra}
\end{equation}
The antiholomorphic modes satisfy the same algebra with central charge
\(\widetilde c\), and
\[
  [L_n,\widetilde L_m]=0.
\]

\begin{figure}[H]
  \centering
  \resizebox{0.92\textwidth}{!}{%
  \begin{tikzpicture}[x=1cm,y=1cm,every node/.style={font=\scriptsize}]
    \tikzset{
      insertion/.style={circle,fill=violet!80!black,inner sep=1.7pt},
      arrow/.style={-{Latex[length=2mm]},line width=0.8pt},
      contour/.style={green!45!black,line width=0.85pt}
    }
    \node[font=\small] at (-3.25,1.95) {nested contours};
    \draw[contour] (-3.25,0) circle (1.15);
    \draw[contour,orange!85!black] (-3.25,0) circle (0.62);
    \node[insertion] at (-3.25,0) {};
    \node[violet!80!black,above right] at (-3.25,0) {\(T(w)\)};
    \node[green!45!black] at (-4.25,1.07) {\(L_n\)};
    \node[orange!85!black] at (-2.52,0.68) {\(L_m\)};

    \draw[arrow] (-1.45,0)--(-0.25,0);
    \node[align=center] at (-0.85,0.48) {use\\\(T T\) OPE};

    \node at (2.45,0.45)
    {\(\displaystyle [L_n,L_m]=(n-m)L_{n+m}
    +\frac{c}{12}(n^3-n)\delta_{n,-m}\)};
    \draw[blue!65!black,line width=0.8pt] (0.1,-0.15)--(4.8,-0.15);
    \node[orange!85!black] at (2.45,-0.62)
    {central term is the fourth-order pole residue};
  \end{tikzpicture}
  }
  \caption{The Virasoro algebra follows by moving stress-tensor contours and
  evaluating the singular part of the \(T T\) OPE.}
  \label{fig:virasoro-tt-contour}
\end{figure}

\section{Plane and Cylinder}

Let \(w=\tau+\ii\sigma\) be a cylinder coordinate with
\(\sigma\sim\sigma+2\pi\), and use
\[
  z=\ee^w.
\]
The stress tensor transforms as
\begin{equation}
  T_{\rm cyl}(w)
  =
  \left(\frac{\dd z}{\dd w}\right)^2T(z)
  +
  \frac c{12}\{z,w\}
  =
  z^2T(z)-\frac c{24},
  \label{eq:stress-plane-cylinder}
\end{equation}
because
\[
  \{z,w\}
  =
  \frac{z'''(w)}{z'(w)}
  -
  \frac32\left(\frac{z''(w)}{z'(w)}\right)^2
  =
  -\frac12.
\]
Thus the cylinder Hamiltonian and momentum are
\begin{equation}
  H_{\rm cyl}
  =
  L_0-\frac c{24}
  +
  \widetilde L_0-\frac{\widetilde c}{24},
  \qquad
  P_{\rm cyl}
  =
  L_0-\frac c{24}
  -
  \left(\widetilde L_0-\frac{\widetilde c}{24}\right).
  \label{eq:cylinder-H-P-virasoro}
\end{equation}
The shift by \(c/24\) is the Casimir energy on the circle.
